% \iffalse meta-comment
%
% hit-report-meta.dtx 是开题中期报告的元信息字段
%
% \fi
%
% \section{开题中期报告的元信息字段}
%
% \changes{v3.2a}{0000/00/00}{元信息字段的声明单独成文件。它说的是「有哪些字段」，
%   与 \file{src/config/} 说的「字段取什么值」是一对，都是数据；
%   字段怎么存、旧名怎么兼容那套机制在 \file{src/utils/hit-meta.dtx}}
% 开题中期报告用到的元信息字段：字段名、旧名映射、\cs{hitsetup} 各族的路由，
% 以及摘要环境。都是清单，不含机制。
%
% ^^A ======================================================================
% ^^A Module report-meta: 元信息字段（\hitsetup 接受的 key）（hit-report）
% ^^A ======================================================================
%    \begin{macrocode}
%<*report-meta>
%% 中文封面

\hit@define@title{zh}
\hit@define@term{report-type-zh}
\hit@define@term{specialty-zh}
\hit@define@term{author-zh}
\hit@define@term{student-id}
\hit@define@term{class-id}
\hit@define@term{affiliation-zh}
\hit@define@term{supervisor-zh}
\hit@define@term{date-zh}
\hit@define@term{enroll-date-zh}

%    \end{macrocode}
%
% \changes{v3.2a}{0000/00/00}{关键词的逗号拆分与 \cs{hitsetup} 的入口挪到
%   \file{src/utils/hit-meta.dtx}}
% 关键词怎么拆、\cs{hitsetup} 的占位定义都见 \file{src/utils/hit-meta.dtx}。
%
% \changes{v3.2a}{0000/00/00}{\cs{hitsetup} 改用 l3keys 做入口，加一层按职责
%   划分的子族路由，与 \file{hithesisbook.cls} 一致。子族暂时只转发，键仍由
%   \texttt{hit/class} 族处理}
% 入口换成 l3keys，根族路由到子族，子族一律把键转回
% \texttt{hit} 族。转发时值要重新加花括号，否则含逗号的值会被二次拆分成多个键。
%
% \changes{v3.2a}{0000/00/00}{元信息、可本地化字符串与参考文献选项正式注册进
%   对应的族，不再依赖子族兜底转发}
% 元信息字段与可本地化字符串。做法与 \file{hithesisbook.cls} 相同：
% 前者的处理器直接调用同名命令，后者的键名由宏名去掉 \texttt{hit@} 前缀、
% 其余 \texttt{@} 换成连字符得来。
%    \begin{macrocode}
\hit@declare@info@field
  {
    title-zh , affiliation-zh , author-zh , class-id , date-zh , enroll-date-zh , student-id ,
    specialty-zh , supervisor-zh , report-type-zh
  }
\hit@declare@constant@as
  {
    equationname , listfigureename , listtableename , listequationename ,
    listequationname , cabstractcname , cabstractename , eabstractcname , eabstractename
  }
\hit@declare@constant@ctex
  {
    chapter-name = chapter/name , appendixname , contentsname , listfigurename ,
    listtablename , figurename , tablename , bibname , indexname
  }
\hit@declare@constant
  {
    campus-full ,
    report-header , report-front-header ,
    bibliography-heading-zh ,
    authorization-usage-heading , author-signature-field-label , date-blank ,
    date-field-label-front , harbin-bachelor-school-bottom-mark , harbin-school-bottom-mark ,
    shenzhen-doctor-midterm-note , shenzhen-master-school-bottom-mark ,
    supervisor-signature-field-label , bachelor-affiliation-field-label ,
    bachelor-affiliation-field-label-shenzhen , bachelor-affiliation-field-label-weihai ,
    bachelor-class-field-label , bachelor-date-field-label ,
    bachelor-check-date-field-label , bachelor-major-field-label ,
    bachelor-student-id-field-label , bachelor-student-field-label ,
    bachelor-supervisor-field-label , bachelor-document-type , bachelor-degree-level ,
    bachelor-teacher-comment-field-label , bachelor-teacher-signature-field-label ,
    affiliation-field-label-zh , associate-supervisor-field-label-zh , author-field-label-zh ,
    co-supervisor-field-label-zh , date-field-label-zh , degree-of-discipline-format-zh ,
    degree-field-label-zh , school-zh , school-field-label-zh , specialty-field-label-zh ,
    supervisor-field-label-zh , thesis-midterm-check , document-type ,
    thesis-title-prefix , degree-level-zh , figure-prefix-en , table-prefix-en ,
    eqdenote-label-zh , eqdenote-label-en , eqdenote-dash , graduate-affiliation-field-label
    , graduate-affiliation-major-field-label , graduate-major-field-label ,
    graduate-date-field-label , graduate-enroll-date-field-label ,
    graduate-student-id-field-label , graduate-student-field-label ,
    graduate-supervisor-field-label , graduate-thesis-field-label , hi ,
    classified-index-international-field-label-zh , classified-index-national-field-label-zh ,
    school-id , school-id-field-label , secret-level-field-label , campus ,
    stage-document-type , stage-midterm , stage-opening , title-separator-zh
  }
%    \end{macrocode}
%
% 参考文献的三个选项，改用连字符命名。
%    \begin{macrocode}
\keys_define:nn { hit / bib }
  {
    label-align .code:n = { \str_set:Nn \g_bib_labelalign_str {#1} } ,
    author-case .code:n =
      { \str_if_eq:nnT {#1} { full } { \bool_set_true:N \g_bib_authorupper_bool } } ,
    max-author .code:n =
      {
        \int_set:Nn \g_bib_maxauthor_int {#1}
        \int_if_zero:nT { \g_bib_maxauthor_int } { \int_gdecr:N \g_bib_maxauthor_int }
      } ,
  }
\keys_define:nn { hit / style }
  {
    ragged-bottom .choice: ,
    ragged-bottom / true  .code:n = { \bool_gset_true:N \g__hit_raggedbottom_bool } ,
    ragged-bottom / false .code:n = { \bool_gset_false:N \g__hit_raggedbottom_bool } ,
    ragged-bottom .default:n = true ,
  }
%    \end{macrocode}
%
% \changes{v3.2a}{0000/00/00}{加旧名到新名的映射表与废弃警告}
% 兼容期的处理，与 \file{hithesisbook.cls} 一致：旧名照常生效，只在日志里提示。
%    \begin{macrocode}
\prop_gset_from_keyval:Nn \g__hit_renamed_option_prop
  {
    raggedbottom   = { style = { ragged-bottom = ... } } ,
    biblabelalign  = { bib = { label-align = ... } } ,
    bibauthorupper = { bib = { author-case = ... } } ,
    bibmaxauthor   = { bib = { max-author = ... } } ,
  }
%    \end{macrocode}
%
% \changes{v3.2a}{0000/00/00}{元信息字段的旧名走废弃期：仍然可用，用一次提示一次}
% 旧名到新名的对照。旧名在 v3.1e 就是公开接口，写在用户的 \file{cover.tex} 里，
% 所以保留到 v4.1。
%    \begin{macrocode}
\hit@define@legacy@terms
  {
    ctype=report-type-zh ,
    csubject=specialty-zh ,
    cauthor=author-zh ,
    cstudentid=student-id ,
    cclassid=class-id ,
    caffil=affiliation-zh ,
    csupervisor=supervisor-zh ,
    cdate=date-zh ,
    cenrolldate=enroll-date-zh
  }
\hit@define@legacy@title{ctitlecover}{zh}{first}
\hit@define@legacy@title{ctitleone}{zh}{first}
\hit@define@legacy@title{ctitletwo}{zh}{second}
\prop_gset_from_keyval:Nn \g__hit_renamed_load_option_prop
  {
    font     = { font-en = ... } ,
    cjk-font = { font-zh = ... } ,
    type     = { level } ,
  }
\__hit_check_class_options:
%    \end{macrocode}
%
% 根族路由。\option{toc} 既是族名又是老的布尔选项名，所以它单独处理：
% 值是 \texttt{true}/\texttt{false} 时当老写法，其余当子族的键值表。
%    \begin{macrocode}
\keys_define:nn { hit }
  {
    info  .meta:nn = { hit / info  } {#1} ,
    style .meta:nn = { hit / style } {#1} ,
    bib   .meta:nn = { hit / bib   } {#1} ,
    const .meta:nn = { hit / const } {#1} ,
    structure .meta:nn = { hit / structure } {#1} ,
    theorem .code:n = { \hit@theorem@style {#1} } ,
    theorem-list .code:n = { \hit@theorem@list {#1} } ,
    theorem-star .code:n = { \hit@theorem@declare@star@list {#1} } ,
    toc .code:n =
      {
        \str_case:nnF {#1}
          {
            { true }  { \bool_gset_true:N \g__hit_toc_bool }
            { false } { \bool_gset_false:N \g__hit_toc_bool }
          }
          { \keys_set:nn { hit / toc } {#1} }
      } ,
    toc .default:n = true ,
    unknown .code:n = { \__hit_legacy_key:n {#1} } ,
  }
\RenewDocumentCommand \hitsetup { +m } { \keys_set:nn { hit } {#1} }
%</report-meta>
%    \end{macrocode}
%
% \changes{v3.2a}{0000/00/00}{\texttt{report-options} 从 \file{src/hit-report.dtx} 挪进来}
% ^^A ======================================================================
% ^^A Module report-options: 文档类选项的声明与后处理（type/campus/fontset 等）（hit-report）
% ^^A ======================================================================
%    \begin{macrocode}
%<*report-options>
%    \end{macrocode}
%
% \changes{v3.2a}{0000/00/00}{加通用选择器 \cs{hit@if@option@TF}，与 \file{hithesisbook.cls}
%   一致}
% 读标志位的那族判断函数（\cs{hit@if@option@TF} 等）定义在
% \file{src/config/hit-options.dtx}，各个类共用。这里只写本类特有的两个：
% \cs{hit@stage@name} 取「开题」或「中期」；按校区、学位、阶段组合出来的判断由
% \cs{hit@define@condition} 生成，一条一行，见下。参数从 expl3 之外传进来，
% 里面的空格照原样保留；替换旧写法时注意 \cs{fi} 后面被跳过的那个空格。
%
% \changes{v3.2a}{0000/00/00}{补上 \option{style/glue} 与 \cs{hit@glue@skip}、
%   \cs{hit@glue@font@size}，与 \file{hithesisbook.cls} 对齐。默认关，版面不变}
% \option{glue} 开时给标题前后距与行距一个伸缩量，关时是固定值。
%    \begin{macrocode}
\cs_new:Npn \hit@stage@name
  {
    \bool_if:NTF \g__hit_opening_bool
      { \hit@stage@opening } { \hit@if@option@T { midterm } { \hit@stage@midterm } }
  }
\hit@define@condition{weihai@or@shenzhen@master}{ ( \g__hit_weihai_bool || \g__hit_shenzhen_bool ) && \g__hit_master_bool }
\hit@define@condition{shenzhen@master@opening}{ \g__hit_shenzhen_bool && \g__hit_master_bool && \g__hit_opening_bool }
\hit@define@condition{harbin@master@opening@or@midterm}{ \g__hit_harbin_bool && \g__hit_master_bool
     && ( \g__hit_opening_bool || \g__hit_midterm_bool ) }
\hit@define@condition{harbin@master@opening}{ \g__hit_harbin_bool && \g__hit_opening_bool && \g__hit_master_bool }
\hit@define@condition{shenzhen@doctor@midterm}{ \g__hit_shenzhen_bool && \g__hit_doctor_bool && \g__hit_midterm_bool }
%    \end{macrocode}
%
% 处理文档类选项
%    \begin{macrocode}

%% 处理文档类选项
%    \end{macrocode}
%
% 键名与变量名不同的字符串选项，\option{cjk-font} 这种带连字符的键要走这个。
%
% \changes{v3.2a}{0000/00/00}{\option{type} 改名 \option{level}，旧名走废弃期。
%   派发类的 \option{kind} 进来之后两个名字都像「种类」，分不清谁管学位谁管交付物}
% \option{level} 是学位层次。原先叫 \option{type}，v3.2a 加了派发类的 \option{kind}
% 之后两个名字都读作「种类」，一个管学位一个管交付物，谁也说不清哪个是哪个。
%
% 旧名 \option{type} 照旧生效，转发到 \option{level}，并提示一次。取值判断看的是
% \cs{g\_\_hit\_}\meta{取值}\cs{\_bool}，跟键名无关，所以内部代码一处不用改。
%    \begin{macrocode}
\clist_map_inline:nn { raggedbottom , toc }
  {
    \hit@declare@boolean@option {#1} { true }
    \bool_new:c { g__hit_#1_given_bool }
    \keys_define:nn { hit / class }
      {
        #1 .code:n =
          {
            \bool_gset_true:c { g__hit_#1_given_bool }
            \str_if_eq:nnTF {##1} { false }
              { \bool_gset_false:c { g__hit_#1_bool } }
              { \bool_gset_true:c { g__hit_#1_bool } }
          } ,
        #1 .default:n = { true } ,
      }
  }
\clist_map_inline:nn { debug , glue , print , newtxmath }
  { \hit@declare@boolean@option {#1} { false } }

\hit@declare@choice@option{level}{}{bachelor,master,doctor}{}
\keys_define:nn { hit / class }
  { type .code:n = { \keys_set:nn { hit / class } { level = #1 } } }

\hit@declare@choice@option{campus}{}{shenzhen,weihai,harbin}{}
\bool_lazy_or:nnF
  { \bool_if_p:N \g__hit_weihai_bool }
  { \bool_if_p:N \g__hit_shenzhen_bool }
  { \bool_gset_true:N \g__hit_harbin_bool }

\hit@declare@choice@option{stage}{}{opening,midterm}{}
%    \end{macrocode}

% \changes{v3.1c}{2023/04/16}{增加用于打印的选项 \meta{print}.\issue[138]}
%    \begin{macrocode}


\hit@declare@string@option{fontset}
%    \end{macrocode}
%
% \changes{v3.2a}{0000/00/00}{西文字体独立成 \option{font} 选项}
% 西文字体，取值与 \file{hithesisbook.cls} 相同。
%    \begin{macrocode}
\hit@declare@string@option@as{font-en}{font@en}
%    \end{macrocode}
%
% \changes{v3.2a}{0000/00/00}{中文字体独立成 \option{cjk-font} 选项}
% 中文字体。留空时照旧把 \option{fontset} 转给 \pkg{ctex}，给了值就由模板自己设置。
%    \begin{macrocode}
\hit@declare@string@option@as{font-zh}{font@zh}
%    \end{macrocode}
%
% \changes{v3.2a}{0000/00/00}{\option{font} 与 \option{cjk-font} 改名
%   \option{font-en} 与 \option{font-zh}，与模板其余选项的中英后缀一致。旧名仍
%   生效，只在日志里提示}
%    \begin{macrocode}
\keys_define:nn { hit / class }
  {
    font     .tl_gset:c = { hit@font@en } ,
    cjk-font .tl_gset:c = { hit@font@zh } ,
  }
%    \end{macrocode}
%
% \changes{v3.2a}{0000/00/00}{新增 \option{math-font} 与 \option{math-style} 选项}
% 数学字体与数学排版风格，取值与 \file{hithesisbook.cls} 相同。
%    \begin{macrocode}
\hit@declare@string@option@as{math-font}{mathfont}
\hit@declare@string@option@as{math-style}{mathstyle}
\keys_define:nn { hit / class }
  {
    unknown .code:n = { \hit@pass@or@warn { ctexart } {#1} } ,
  }
\PassOptionsToPackage{no-math}{fontspec}

\ProcessKeyOptions [~hit~/~class~]

%    \end{macrocode}
%
% \changes{v3.2a}{0000/00/00}{art-options 的三段选项校验改用 \cs{bool\_lazy\_any:nF}、
%   \cs{bool\_lazy\_or:nnF} 与 \cs{bool\_lazy\_all:nT}，嵌套的 \cs{relax}\cs{else}
%   链换成显式判断。报错文本含英文与空格，按 CONTRIBUTING §4.7 留在 expl3 之外}
% 三段报错文本原样留在 expl3 之外，里面的空格与花括号都有意义。
%    \begin{macrocode}
\cs_new_protected:Npn \hit@report@error@type
  {
    \ClassError { hit-report }
      {
        \MessageBreak Please~specify~thesis~type~in~option:~
        \MessageBreak type=[bachelor~|~master~|~doctor]~
      }
  }
\cs_new_protected:Npn \hit@report@error@stage
  {
    \ClassError { hit-report }
      {
        \MessageBreak Please~specify~stage~in~option:~
        \MessageBreak stage=<opening|midterm>~
      }
  }
\bool_lazy_any:nF
  {
    { \bool_if_p:N \g__hit_bachelor_bool }
    { \bool_if_p:N \g__hit_master_bool }
    { \bool_if_p:N \g__hit_doctor_bool }
  }
  { \hit@report@error@type }
\bool_lazy_or:nnF
  { \bool_if_p:N \g__hit_opening_bool }
  { \bool_if_p:N \g__hit_midterm_bool }
  { \hit@report@error@stage }
\hit@if@shenzhen@doctor@midterm
  {
    \bool_if:NF \g__hit_toc_given_bool { \bool_gset_false:N \g__hit_toc_bool }
    \bool_if:NF \g__hit_raggedbottom_given_bool
      { \bool_gset_false:N \g__hit_raggedbottom_bool }
  }
%    \end{macrocode}
%
% \changes{v3.2a}{0000/00/00}{选项后处理（AutoFakeBold 与 fontset 透传）从 artcls 移入 art-options 模块}
%    \begin{macrocode}
\RequirePackage{ifthen}
\tl_if_empty:NTF \hit@font@zh
  {
    \PassOptionsToPackage { AutoFakeBold=2 } { xeCJK }
    \tl_if_empty:NF \hit@fontset
      { \PassOptionsToClass { fontset=\hit@fontset } { ctexart } }
  }
  { \PassOptionsToClass { fontset=none } { ctexart } }
%</report-options>
%    \end{macrocode}
%
% \Finale
